\documentclass[11pt,a4paper]{article}

\usepackage[utf8]{inputenc}
\usepackage[T1]{fontenc}
\usepackage[french]{babel}
\usepackage{geometry}
\usepackage{xcolor}
\usepackage{hyperref}
\usepackage{tabularx}
\usepackage{longtable}
\usepackage{enumitem}
\usepackage{listings}
\usepackage{titlesec}

\geometry{margin=2.4cm}
\hypersetup{
  colorlinks=true,
  linkcolor=blue!55!black,
  urlcolor=blue!55!black
}

\definecolor{chorusblue}{HTML}{1D4ED8}
\definecolor{chorusgray}{HTML}{F3F4F6}
\definecolor{chorusdark}{HTML}{111827}

\titleformat{\section}
  {\Large\bfseries\color{chorusblue}}
  {\thesection}{0.75em}{}

\titleformat{\subsection}
  {\large\bfseries\color{chorusdark}}
  {\thesubsection}{0.75em}{}

\lstdefinestyle{chorus}{
  backgroundcolor=\color{chorusgray},
  basicstyle=\ttfamily\small,
  breaklines=true,
  frame=single,
  rulecolor=\color{black!15},
  columns=fullflexible
}
\lstset{style=chorus}

\setlist[itemize]{leftmargin=1.2cm}
\setlist[enumerate]{leftmargin=1.2cm}

\title{
  \vspace{-1.5cm}
  \textbf{\Huge Automatisation GitHub}\\
  \vspace{0.3cm}
  \Large Workflows CI/CD du projet Chorus Pilot
}
\author{Projet Chorus}
\date{\today}

\begin{document}

\maketitle
\thispagestyle{empty}

\begin{center}
  \vspace{0.5cm}
  \begin{tabularx}{0.9\textwidth}{|X|}
    \hline
    \textbf{Résumé} \\
    Le projet Chorus Pilot utilise GitHub Actions pour automatiser les
    contrôles qualité, les tests, la sécurité, la validation Docker et le
    déploiement de l'environnement de développement. Cette automatisation permet
    de sécuriser les changements avant leur intégration et de réduire les actions
    manuelles lors des mises à jour. \\
    \hline
  \end{tabularx}
\end{center}

\newpage
\tableofcontents
\newpage

\section{Objectif de l'automatisation}

Le dossier \texttt{.github/workflows} contient les scénarios GitHub Actions du
projet. Ces fichiers décrivent ce qui doit être exécuté automatiquement lorsqu'un
développeur pousse du code, ouvre une pull request ou lance manuellement un
workflow.

L'objectif est de répondre à quatre besoins :

\begin{itemize}
  \item vérifier que le frontend fonctionne toujours ;
  \item vérifier que le backend Python reste valide ;
  \item contrôler les risques de sécurité ;
  \item déployer automatiquement la branche de développement quand les contrôles
        sont réussis.
\end{itemize}

\section{Branches surveillées}

Les workflows principaux sont déclenchés sur les branches importantes du projet :

\begin{itemize}
  \item \texttt{main} ;
  \item \texttt{PRE\_PROD\_V0} ;
  \item \texttt{developpement}.
\end{itemize}

Ils se lancent lors d'un \textit{push}, d'une \textit{pull request} ou d'un
lancement manuel depuis l'interface GitHub Actions.

\section{Workflow CI}

Le fichier \texttt{.github/workflows/ci.yml} met en place la chaîne de contrôle
principale du projet. Il est nommé \textbf{CI}.

\subsection{Contrôles frontend}

Le job frontend exécute les étapes suivantes :

\begin{enumerate}
  \item récupération du code du dépôt ;
  \item installation de Node.js 20 ;
  \item installation propre des dépendances avec \texttt{npm ci} ;
  \item vérification du code avec \texttt{npm run lint} ;
  \item exécution des tests avec \texttt{npm run test} ;
  \item génération du build avec \texttt{npm run build}.
\end{enumerate}

Ce job permet de vérifier que l'interface React compile correctement, que les
tests passent et que le code respecte les règles de qualité définies dans le
projet.

\subsection{Contrôles backend}

Le job backend vérifie la partie Python :

\begin{enumerate}
  \item installation de Python 3.12 ;
  \item installation des dépendances depuis \texttt{requirements.txt} ;
  \item compilation des modules Python avec \texttt{compileall} ;
  \item exécution des tests Python avec \texttt{unittest}.
\end{enumerate}

Ce contrôle évite d'intégrer du code Python qui ne compile pas ou qui casse les
tests existants.

\subsection{Validation Docker Compose}

Un troisième job valide la configuration Docker de production :

\begin{enumerate}
  \item copie du fichier \texttt{.env.example} vers \texttt{.env.local} pour
        fournir une configuration minimale ;
  \item validation du fichier \texttt{docker-compose.prod.yml} avec
        \texttt{docker compose config}.
\end{enumerate}

Cette étape vérifie que la stack Docker reste cohérente et que le fichier Compose
peut être interprété correctement.

\section{Workflow Security}

Le fichier \texttt{.github/workflows/security.yml} ajoute des contrôles de
sécurité. Il est nommé \textbf{Security}.

\subsection{Protection contre les secrets suivis par Git}

Le premier contrôle bloque la présence de fichiers sensibles dans le dépôt :

\begin{itemize}
  \item fichiers \texttt{.env} ou dérivés ;
  \item clés privées \texttt{.pem} ;
  \item fichiers \texttt{.key}.
\end{itemize}

Le fichier \texttt{.env.example} reste autorisé, car il sert de modèle sans
contenir de vrais secrets.

\subsection{Recherche de secrets évidents}

Le workflow recherche aussi des variables sensibles qui auraient été commitées
par erreur, par exemple :

\begin{itemize}
  \item secret JWT applicatif ;
  \item hash de mot de passe applicatif ;
  \item secret client Chorus ;
  \item mot de passe technique Chorus ;
  \item secrets exposés dans des variables frontend.
\end{itemize}

Si une valeur suspecte est trouvée, le workflow échoue pour empêcher l'intégration
du changement.

\subsection{Audit des dépendances frontend}

Le job \texttt{npm-audit} installe les dépendances frontend et exécute :

\begin{lstlisting}
npm audit --audit-level=high
\end{lstlisting}

L'objectif est de détecter les vulnérabilités importantes dans les dépendances
JavaScript.

\subsection{Audit des dépendances Python}

Le job \texttt{pip-audit} installe l'outil \texttt{pip-audit} puis analyse les
dépendances listées dans \texttt{requirements.txt}. Il permet d'identifier les
vulnérabilités connues côté backend Python.

\subsection{Exécution planifiée}

En plus des push et pull requests, le workflow sécurité est lancé automatiquement
chaque semaine. Cela permet de détecter une vulnérabilité même si le code n'a pas
changé récemment.

\section{Workflow Deploy Development}

Le fichier \texttt{.github/workflows/deploy-development.yml} automatise le
déploiement sur le serveur de développement. Il est nommé
\textbf{Deploy Development}.

\subsection{Condition de déclenchement}

Le déploiement se lance automatiquement lorsque le workflow \textbf{CI} se
termine avec succès sur la branche \texttt{developpement}. Il peut aussi être
lancé manuellement depuis GitHub Actions.

Cette logique évite de déployer une version qui n'a pas passé les contrôles de
base.

\subsection{Secrets nécessaires}

Le workflow utilise des secrets GitHub pour se connecter au serveur :

\begin{itemize}
  \item \texttt{DEV\_SSH\_HOST} : adresse du serveur ;
  \item \texttt{DEV\_SSH\_USER} : utilisateur SSH ;
  \item \texttt{DEV\_SSH\_PRIVATE\_KEY} : clé privée SSH.
\end{itemize}

Ces informations ne sont pas stockées dans le dépôt. Elles sont configurées dans
les paramètres GitHub du projet.

\subsection{Variables optionnelles}

Deux variables peuvent personnaliser le déploiement :

\begin{itemize}
  \item \texttt{DEV\_APP\_DIR} : dossier de l'application sur le serveur ;
  \item \texttt{DEV\_COMPOSE\_PROJECT} : nom du projet Docker Compose.
\end{itemize}

Si elles ne sont pas définies, le workflow utilise des valeurs par défaut.

\subsection{Étapes réalisées sur le serveur}

Une fois connecté en SSH, le workflow :

\begin{enumerate}
  \item se place dans le dossier de l'application ;
  \item récupère les dernières modifications de la branche \texttt{developpement} ;
  \item bascule sur cette branche ;
  \item applique un \texttt{git pull --ff-only} ;
  \item reconstruit et relance la stack Docker ;
  \item affiche l'état des conteneurs.
\end{enumerate}

La commande principale de redéploiement est :

\begin{lstlisting}
docker compose -p "$DEV_COMPOSE_PROJECT" \
  --env-file .env.local \
  -f docker-compose.prod.yml \
  up -d --build
\end{lstlisting}

\section{Dependabot}

Le fichier \texttt{.github/dependabot.yml} configure Dependabot pour surveiller
les mises à jour de dépendances.

Dependabot vérifie chaque semaine :

\begin{itemize}
  \item les dépendances npm du frontend ;
  \item les dépendances pip du backend Python ;
  \item les actions GitHub utilisées par les workflows.
\end{itemize}

Le nombre de pull requests ouvertes automatiquement est limité à cinq par
écosystème afin de garder un volume de maintenance raisonnable.

\section{Résultat obtenu}

Grâce à cette configuration, le projet dispose d'un premier pipeline CI/CD
complet :

\begin{itemize}
  \item les erreurs frontend sont détectées avant intégration ;
  \item les erreurs backend sont détectées avant intégration ;
  \item la configuration Docker est validée automatiquement ;
  \item les secrets commis par erreur sont bloqués ;
  \item les dépendances sont auditées régulièrement ;
  \item les mises à jour de dépendances sont proposées automatiquement ;
  \item l'environnement de développement peut être redéployé automatiquement
        après une CI réussie.
\end{itemize}

\section{Conclusion}

Le travail réalisé côté \texttt{.github/workflows} apporte une automatisation
essentielle au projet Chorus Pilot. Il sécurise le cycle de développement,
réduit les vérifications manuelles et prépare une base solide pour faire évoluer
le projet vers une chaîne d'intégration et de déploiement plus complète.

\end{document}
